\begin{table*}[t] \centering \caption{Positioning against prior art.} \label{tab:rw_compare} \begin{tabular}{llll} \toprule \textbf{Category} & \textbf{Representative} & \textbf{Key Traits} & \textbf{Our Distinction} \\ \midrule Fairness-centric & CFS/ULE & Strong fairness, higher tails & Bounded bias with fairness floors \\ Low-latency desktop & BFS/MuQSS & Interactive focus, weaker fairness & Kernel-native bias + guardrails \\ QoS feedback & Quasar-like & SLO-driven, userspace hooks & In-kernel, minimal operator surface \\ ML-guided & iTune, TerrierTail & Predictive, opaque failure modes & Simple, explainable hybrid control \\ \bottomrule \end{tabular} \end{table*}

\section{Related Work}
\label{sec:related}
We position BRS against the scheduling and QoS literature summarized in Table~\ref{tab:rw_compare}, spanning fairness-centric, low-latency, and ML-based schedulers in general-purpose, cloud, and embedded settings.

\subsection{Evolution of Fairness--Latency Trade-offs in Kernel Scheduling}
Mainstream schedulers such as Linux CFS and FreeBSD ULE rely on proportional allocation, prioritizing long-term fairness and starvation prevention~\cite{241774,4624248,67577}. In practice, however, strict fairness can be counterproductive: under contention, latency spikes for the interactive tasks users care about most~\cite{4663063,6365626}. BFS and MuQSS target this by favoring responsiveness~\cite{10306992}, and prior work suggests ``overly fair'' scheduling can harm interactive workloads~\cite{Bouron2018Battle}. The same tension appears in heterogeneous multicore research, where throughput and latency oppose each other~\cite{8360522,6899614}. BRS builds on the observation that even minor changes to work accounting can substantially improve responsiveness, favoring precise, limited kernel changes over complex overhauls.

\subsection{Classical Interactivity Heuristics and Mainline Mechanisms}
\label{subsec:rw-mlfq}
Using observed sleep and I/O behavior to favor interactive tasks is a long-standing idea, and a fair positioning of BRS must make the relationship explicit. \emph{Multilevel Feedback Queue} (MLFQ) schedulers maintain several priority queues and \emph{promote} a task to a higher-priority queue when it is observed to sleep or block, demoting it when it is CPU-bound. The \emph{Linux O(1) scheduler}, which CFS replaced in 2007, used essentially the same principle: it granted an interactivity bonus computed from a task's \texttt{sleep\_avg} and adjusted the task's effective priority accordingly. BRS shares the input signals of these designs---sleep duration and I/O behavior---but differs in three concrete ways. (i)~\emph{Mechanism:} BRS expresses the preference as a continuous, \emph{explicitly bounded} offset on \textit{vruntime} (Definition~\ref{def:bound}) rather than a discrete queue-level promotion, so it composes with CFS's proportional-fairness accounting instead of replacing it. (ii)~\emph{Guarantees:} classical heuristics offer no formal stability or fairness guarantee, whereas BRS provides the bound of Lemma~\ref{lem:fairness} and the mean-square-stability result of Theorem~1. (iii)~\emph{Gaming resistance:} MLFQ-style promotion is well known to be exploitable by tasks that sleep strategically just before their quantum expires; BRS's bounded, aged score and controller back-pressure (Section~\ref{subsec:gaming}) limit this leverage by construction.

Mainline Linux has also moved toward exposing responsiveness as a tunable rather than a fixed policy. The \emph{utilization clamping} (UClamp) interface lets user space hint a minimum/maximum utilization per task or cgroup, primarily to steer frequency selection and big.LITTLE placement, and the \emph{latency\_nice} attribute biases wake-up preemption decisions for latency-sensitive tasks. These mechanisms are closely related to BRS's motivation, but they act through per-task hints supplied from user space and do not bound or formally analyze the resulting fairness deviation. BRS is complementary: it keeps the control loop inside the kernel, requires no per-application hint, and adds an explicit bound and a stability argument. A scheduler could in principle consume UClamp/latency\_nice hints as an additional input to the interactivity score, which we leave as future work.

\subsection{Predictive and Responsiveness-Oriented Scheduling}
A long line of work has explored letting the kernel \emph{anticipate} task behavior rather than react to it. Early micro-schedulers used short-burst prediction to improve responsiveness~\cite{4663063}; later designs incorporated thermal signals into scheduling decisions~\cite{6365626}. While promising in the literature, our experience is that unreliable predictions in production are a liability rather than an asset. As a result, BRS deliberately avoids opaque black-box models and instead uses a strictly constrained \textit{vruntime} bias. This is a pragmatic choice: it improves user-perceived latency without violating ordering semantics. For real-time weakly hard systems~\cite{9351534}, we likewise find that regular, predictable structure matters more than chasing subtle but unstable performance gains.

\subsection{QoS-Aware and Feedback-Controlled Scheduling}
QoS-aware scheduling balances application needs against resource limits, often via feedback loops; cloud-native tools such as iTune push this with heuristics~\cite{5580094,5169938}. Their recurring weakness is heavy reliance on user space, which adds complexity and production-risky failure modes~\cite{7792746,6782279,8758850}. BRS instead embeds QoS-relevant bias directly in the kernel with locked-in fairness floors, replacing brittle black-box controllers with a minimal, interpretable surface that cuts overhead and opacity~\cite{7426385,8338088}.

\subsection{Mobile, Embedded, and Edge Scheduling Contexts}
Mobile and edge devices are fundamentally constrained by power and thermal limits~\cite{1244943}. While Energy-Aware Scheduling (EAS) and jitter-aware policies~\cite{9121657} target smoothness, IoT operating systems frequently sacrifice flexibility to preserve determinism~\cite{10007654,9999639}, an undesirable but often unavoidable trade-off. Even ML-driven tail-latency methods~\cite{Zhang2025Tail} do not provide what these settings most require: a stable, low-level scheduling foundation.

A complementary body of work targets energy- and reliability-aware scheduling on heterogeneous platforms. HEAT couples deadline partitioning and core clustering with temperature- and energy-aware scheduling on DVFS-enabled cores, improving task-set acceptance while lowering temperature~\cite{Sharma2022HEAT}; FRESH adds fault tolerance via a semi-partitioned, job-progress-monitoring scheme with replicated copies and selective processor sleep~\cite{Moulik2024FRESH}. These optimize system-level objectives---energy, temperature, reliability---largely orthogonal to BRS's tail-latency objective, but a deployable scheduler for heterogeneous mobile SoCs must eventually reconcile responsiveness with them. BRS is positioned as the responsiveness-and-fairness substrate on which such energy/thermal/reliability policies can sit, and Section~\ref{subsec:future} discusses integrating BRS with a DVFS governor in this spirit.

\subsection{Positioning of BRS}
Prior approaches fall broadly into three categories: (i)~classical fairness-centric schedulers and the proportional-allocation theory behind them~\cite{241774,4624248}; (ii)~low-latency, interactivity-first kernel schedulers such as BFS and MuQSS and the classical MLFQ/O(1) heuristics of Section~\ref{subsec:rw-mlfq}, which favor responsiveness but weaken fairness~\cite{10306992,4663063}; and (iii)~QoS- and ML-driven frameworks that adapt to workload signals but rely on user-space hooks and opaque models~\cite{7426385,8338088,Zhang2025Tail}. BRS occupies the middle ground: rather than a major redesign, it leverages an explicit fairness lower bound to impose a limited, bounded response bias on the kernel's hot path, so that no task is starved. By complementing---rather than replacing---higher-level ML or QoS policies, BRS enables practical configurations across real-world hardware.